% 张三慧第4版第二卷第23章《光的衍射》原书思考题23.1--23.9
% V2 PDF物理页287（印刷270）；与“习题23.1--23.19”分开计数。
\section*{第23章\quad 光的衍射：思考题}
\paragraph{思考题 23.1}
为什么日常生活中声波绕过障碍物的衍射现象，
通常比可见光绕过类似障碍物的衍射更容易观察？
\paragraph{思考题 23.2}
在观察单缝夫琅禾费衍射的装置中，
缝后会聚透镜起什么作用？其焦平面为何能呈现远场角分布？
\paragraph{思考题 23.3}
某观察方向对应单缝波阵面分成四个等面积半波带
（图23.30）；第一与第三带相位相同，
第二与第四带相位相同。
为什么屏幕上的P点仍为衍射\emph{极小}
而不是极大？
\paragraph{思考题 23.4}
在用透镜观察单缝夫琅禾费衍射时，
(1) 把单缝沿垂直于透镜光轴方向平移，
焦平面衍射图样是否变化？
(2) 若将有限距离的狭缝光源沿垂直光轴方向平移，
使入射倾角改变，图样又会怎样变化？
\paragraph{思考题 23.5}
单缝夫琅禾费衍射中逐渐加宽缝宽，
中央明纹的角宽度、方向分辨率和
通过狭缝的总光通量分别如何变化？
\paragraph{思考题 23.6}
假设人眼能感光的电磁波从可见波段变成
厘米波段，但瞳孔直径仍只有约3mm。
人还能像现在这样分辨外部物体轮廓吗？
请用衍射角分辨尺度判断。
\paragraph{思考题 23.7}
保持多缝光栅的缝距 \(d\) 不变，
而改变等间距缝的总数 \(N\)。
各级主极大的理想角位置会否变化？
峰宽、峰强与缺级现象如何变化？
\paragraph{思考题 23.8}
对杨氏双缝，
(1) 只开第一缝；
(2) 只开第二缝；
(3) 两缝同时打开。
描述三种强度分布，
并说明为何前两个强度分布之和不等于第三种。
\paragraph{思考题 23.9}
许多相同缝宽的狭缝沿直线排布，
但相邻缝距大小随机且缺乏周期性。
单色平行光正入射时，
衍射图样是否还会出现规则、锐利的
非零级光栅主极大？中央方向又如何？
